Work in the \(r=2nl-l+1\) free coordinates of \(\mathcal G_{n,l}\), inserting the \(l-1\) fixed zeros when a type row is evaluated. First validate the explicit type encodings in \(B\) and remove duplicate list entries, since the signature is a set. For every \(\tau=(\delta,\eta)\in B\), form its zero--one normal \(\alpha_\tau\) in the free coordinates and the right-hand side \((c_{n,l})_\tau=\eta(1)-\eta(0)\). Apply exact elimination to
\[\alpha_\tau^{\mathsf T}x=(c_{n,l})_\tau,\qquad \tau\in B.\]
Reject if the system is inconsistent or its row rank is smaller than \(r\).

We justify this rank test without assuming full dimensionality. Set every one of the \(2n\) coordinates in arm one equal to \(n\), set the gauged coordinates \(\lambda_{(0,0,z)}\) for \(z\ge2\) equal to zero, and set every other free coordinate to a positive number. This vector belongs to \(\mathcal G_{n,l}\). Each type sum receives \(n\) from arm one and nonnegative contributions from all other arms, whereas \(\eta(1)-\eta(0)\le n-1\). Every defining inequality is therefore strict. Thus \(\mathcal H^0_{n,l}\) is full-dimensional in its \(r\)-dimensional gauge space for every \(l\ge2\). At a feasible point, active normals span that space if and only if the point is extreme: spanning makes the common active affine set a singleton, while failure to span supplies a nonzero direction orthogonal to every active normal, and a sufficiently small displacement in both signs preserves all inactive strict inequalities.

If the equality system has rank \(r\), exact elimination produces its unique rational candidate \(\lambda\). For each \((a,b)\in\mathcal Y_n^2\), define
\[m_{ab}(\lambda)=\sum_{z=1}^l\min\{\lambda_{(a,0,z)},\lambda_{(b,1,z)}\}.\]
For a type with \(\eta(0)=a\) and \(\eta(1)=b\), minimization over the treatment word separates arm by arm, giving the algebraic identity
\[\min_{\delta\in\mathcal D^{\mathcal Z_l}}\left\{\sum_{z=1}^l\lambda_{(\eta(\delta(z)),\delta(z),z)}-(b-a)\right\}=m_{ab}(\lambda)-(b-a).\]
Consequently \(\lambda\in\mathcal H^0_{n,l}\) if and only if \(m_{ab}(\lambda)\ge b-a\) for all \(n^2\) pairs. Reject if any of these inequalities fails.

Suppose feasibility holds. If \(m_{ab}(\lambda)>b-a\), no type with outcome pair \((a,b)\) is tight. If equality holds, set
\[R_{ab,z}=\operatorname*{argmin}_{d\in\{0,1\}}\lambda_{(\eta(d),d,z)}.\]
For every treatment word, its slack above the separable minimum is a sum of nonnegative armwise gaps. It vanishes exactly when \(\delta(z)\in R_{ab,z}\) for every arm. Hence the tight words for \((a,b)\) form the Cartesian product \(\prod_z R_{ab,z}\), with cardinality
\[2^{t_{ab}},\qquad t_{ab}=|\{z:\lambda_{(a,0,z)}=\lambda_{(b,1,z)}\}|.\]
For each outcome pair, check that every listed word belongs to this product and that the number of distinct listed words is zero in the strict case and \(2^{t_{ab}}\) in the equality case. Inclusion together with equality of finite cardinalities is equivalent to having listed the entire product. Therefore the algorithm accepts exactly when \(B=\operatorname{Sig}(\lambda)\). An accepted candidate is feasible and has \(r\) independent active normals, so it is a vertex. Conversely, the complete signature of any vertex is consistent, has rank \(r\), passes feasibility, and passes every product test. This proves soundness and completeness.

It remains to bound bit complexity. There are \(r=O(nl)\) variables, \(|B|\) explicitly supplied zero--one rows, and \(O(n^2l)\) rational comparisons in the separable tests. Any nonsingular \(r\)-row subsystem has determinant of absolute value at most \(r!\) by expanding the determinant into \(r!\) signed products of zero--one entries. Replacing one column by a right-hand side of magnitude at most \(n-1\) gives the analogous bound \(n r!\). Cramer's rule therefore bounds every numerator and denominator by integers of \(O(r\log r+\log n)\) bits. Fraction-free elimination uses polynomially many arithmetic operations on integers of polynomial bit length. Finally, \(2^{t_{ab}}\) has at most \(l+1\) bits and is compared only with an explicitly counted sublist of \(B\). All validation, elimination, feasibility, completeness, and rank operations thus take bit-operation time polynomial in \(n+l+\operatorname{enc}(B)\), for every \(n,l\ge2\).